\section{Related Works}

\subsection{LLM Post-Training: SFT and RL}

Current post-training methodologies for LLMs are largely centered around two primary paradigms: Supervised Fine-Tuning (SFT) and Reinforcement Learning (RL) \citep{wei2021finetuned, ouyang2022training}. In the SFT paradigm, models are adapted for specific applications through training on curated input-output pairs, a process which has been shown to effectively align their behavior with human demonstrations \citep{chung2022scaling, longpre2023flan,touvron2023llama, touvron2023llama2}. 
In parallel, numerous works have highlighted RL as an effective approach for refining LLM behavior in ways that are difficult to capture with SFT's static datasets \citep{glaese2022improving, bai2022training, nakano2021webgpt}. Within this domain, a popular framework is Reinforcement Learning from Human Feedback (RLHF), which optimizes the LLM policy against a reward model trained on human preferences \citep{christiano2017deep, stiennon2020learning}. Multiple works have established Proximal Policy Optimization (PPO) as a cornerstone algorithm for this phase \citep{schulman2017proximal, ziegler2019fine}. To further improve reasoning capabilities in reward-driven optimization, recent advancements like Group Relative Policy Optimization (GRPO) have also been developed and widely adopted \citep{shao2024deepseekmath, zheng2025group, chen2025minimax}.


\subsection{A Combination of Online and Offline Data in LLM Post-Training}


Beyond applying SFT or RL in isolation, further explorations have sought to synergize their respective strengths by combining signals from pre-existing \textit{offline data} and dynamically generated \textit{online data} \citep{fu2025srft, yan2025learning}. This motivation stems from the distinct characteristics of each approach: SFT is noted for its efficiency in distilling knowledge from offline sources, whereas RL is valued for fostering exploration through online rollouts, a process frequently linked to improved generalization \citep{rajani2025scalpel, chu2025sftmemorizesrlgeneralizes}. The strategies for this integration are diverse; some techniques use offline data as a prefix to guide online generation \citep{zhou2023prefix, touvron2024context, li2025self_guided, wang2025beyond}, while others enhance offline data by incorporating reward signals in a process known as reward-augmented fine-tuning \citep{liu2024dft, zhao2023reward, park2025reward_reweighted, sun2025ktae}. The broader landscape also includes various purely offline preference optimization methods, though they follow a different paradigm \citep{rafailov2023dpo, mitchell2024leveraging, liu2025principled, ethayarajh2024kto, ahmadian2024back}. However, the most direct approach to synergy involves the concurrent use of both data types for training updates.


This direct approach, often termed mix-policy learning, is particularly relevant to our work and typically involves updating the model with a composite objective that combines an SFT loss from offline data and an RL loss from online data \citep{dong2023raft, gulcehre2023reinforced, singh2023beyond, liu2023bridging}. For instance, LUFFY \citep{yan2025learning} explores this paradigm by combining a fixed ratio of offline demonstration data with online rollouts in each training batch. Subsequently, SRFT \citep{fu2025srft} proposed a monolithic training phase that dynamically adjusts the weights of SFT and RL losses based on the model's policy entropy, further demonstrating the viability of unifying these signals over a sequential pipeline. The principle of creating such a composite loss is shared by a variety of other recent frameworks \citep{wu2024alternating, zhang2025on, kim2025dynamic, yu2025dapo, 2025empoweringsmallvlmsthink}, and AMFT \citep{2025aligningllm} begins to explore meta-gradient-based controllers. While these methods highlight a clear trend towards unifying training signals, a foundational theoretical analysis explaining why these different learning signals can be effectively combined is still lacking. This motivates our work to establish a unified theoretical framework that in turn inspires a more principled algorithm design.

